%! Graphing Exponential Functions and Transformations — Unit 6, Lesson 6.2 — Warm-Up (Answer Key)
\documentclass[10pt]{article}
\usepackage{algebra2-article}
\usepackage{algebra2-key}   % pulls in -boxes; do NOT also load -boxes
\newcommand{\TallMath}[1]{$\displaystyle #1\rule[-1.4em]{0pt}{3.2em}$}

\begin{document}

\pageheader{Unit 6, Lesson 6.2}{Warm-Up --- Answer Key}
\namedateperiod

\vspace{0.10in}

\begin{notesbox}{Getting Ready: Skills We'll Use Today}
\small Today you move an exponential graph around. These three skills are the tools. Work quickly.

\vspace{4pt}
\textbf{1. Shifts you already know} (Units 2 and 5). The parent graphs are $y=x^{2}$ and
$y=\sqrt{x}$. Fill in the direction and the amount.

\vspace{3pt}
\hspace*{1.2em} $y=(x-3)^{2}+1$ \; moves \ans{right} $3$ \; and \ans{up} $1$. \\[3pt]
\hspace*{1.2em} $y=\sqrt{x+2}-5$ \; moves \ans{left} $2$ \; and \ans{down} $5$. \\[3pt]
\hspace*{1.2em} A number \emph{outside} the function moves the graph \ans{up or down}; a number
\emph{inside}, with the $x$, moves it \ans{left or right}.

\vspace{6pt}
\textbf{2. Reading a model (Lessons 6.0--6.1).} Everything below comes from $y=3(2)^{x}$ alone.

\vspace{3pt}
\hspace*{1.2em} $a=$ \ans{3} \quad $b=$ \ans{2} \quad Growth or decay?
\ans{growth} \\[4pt]
\hspace*{1.2em} $y$-intercept: \ans{$(0,3)$} \quad Horizontal asymptote: \ans{$y=0$} \quad
Range: \ans{$(0,\infty)$}

\vspace{6pt}
\textbf{3. A negative exponent, two ways} (Unit 5 machinery). Evaluate each one.

\vspace{3pt}
\hspace*{1.2em} $2^{-2}=$ \ans{$\frac14$} \qquad $\left(\dfrac{1}{2}\right)^{2}=$ \ans{$\frac14$}
\qquad $2^{-3}=$ \ans{$\frac18$} \qquad $\left(\dfrac{1}{2}\right)^{3}=$ \ans{$\frac18$}

\vspace{4pt}
\hspace*{1.2em} So are $y=2^{-x}$ and $y=\left(\dfrac{1}{2}\right)^{x}$ two different functions, or
the same function written two ways? \ans{the same function}

\end{notesbox}

\begin{teachernote}
\textbf{Item 1 is the sign flip, and it is the only part of today that is not new.} The rule students
need out loud is the one on the third line: \emph{outside moves it vertically, inside moves it
horizontally, and the inside one goes the ``wrong'' way.} Today's twist is that for an exponential
these two are no longer symmetric --- only the outside one moves the asymptote. Do not say that yet;
the Hook earns it.

\textbf{Item 2} is a straight recall check. Anyone who cannot produce $y=0$ and $(0,\infty)$ here
will not be able to tell what changed after a shift, because the whole lesson is measured against
those two answers. Note the range: the parenthesis at $0$ is Lesson 6.0's headline and it comes back
today as $(k,\infty)$.

\textbf{Item 3 earns the reflection rule numerically before it is stated.} Both columns give the same
numbers, so $2^{-x}$ and $\left(\frac12\right)^{x}$ must be the same function --- which is why
reflecting a growth curve across the $y$-axis produces \emph{decay}. This is not a new rule; it is
the negative-exponent property from Unit 5. Write $2^{-x}=\dfrac{1}{2^{x}}=\left(\dfrac12\right)^{x}$
on the board and leave it up --- Section 3 of the notes points back at it.
\end{teachernote}

\end{document}
